% TOP_PROBLEM: {"problemId": "problem.whitehead-asphericity-conjecture", "canonicalTitle": "Whitehead Asphericity Conjecture", "releaseRank": 397, "primaryDomain": "geometry_topology", "primaryDomainLabel": "Geometry & topology", "displayStatus": "open_disputed_claim", "releaseStatus": "open_disputed_claim"}
% !TeX program = lualatex
% Definition and sources only; this file is not an LLM solution attempt.
\documentclass[11pt]{article}
\usepackage[margin=25mm]{geometry}
\usepackage{fontspec}
\setmainfont{Latin Modern Roman}
\usepackage{amsmath,amssymb,mathtools}
\usepackage{unicode-math}
\setmathfont{Latin Modern Math}
\usepackage{xurl,hyperref}
\hypersetup{colorlinks=true,urlcolor=blue}
\setcounter{secnumdepth}{0}
\setlength{\parindent}{0pt}
\setlength{\parskip}{5pt}
\setlength{\emergencystretch}{3em}
\begin{document}
\section*{\texorpdfstring{Whitehead Asphericity Conjecture}{Whitehead Asphericity Conjecture}}\label{397}
\subsection{Definitions and mathematical statement}
A connected CW-complex $X$ is aspherical if $\pi_i(X)=0$ for all $i\ge2$, equivalently its universal cover is contractible. A subcomplex is a union of cells closed under taking attaching boundaries. Whitehead's conjecture asks
\[
 X\text{ aspherical},\quad\dim X\le2,\quad
 Y\subseteq X\text{ a connected subcomplex}
 \quad\Longrightarrow\quad Y\text{ aspherical}.
\]
Neither finiteness of the complexes nor a special fundamental group is assumed in the recorded scope. In presentation language, a presentation complex has one two-cell per relator; removing relators gives the corresponding subpresentation version. Merely knowing that the inclusion induces an injection on fundamental groups would not by itself be a definition of asphericity.
\par\smallskip\textit{Formulation and concept references:} \href{https://www.mathnet.ru/php/archive.phtml?jrnid=im&option_lang=rus&paperid=9597&wshow=paper}{[S1]}.

\subsection{Short English statement}
Must every connected subcomplex of an aspherical two-dimensional cell complex also have contractible universal cover?
\subsection{Sources}
\begin{itemize}
\item[S1]A. M. Mikhovich, Rational and p-adic analogues of J. H. C. Whitehead's conjecture, Izvestiya: Mathematics 89:2 (2025), section 1.
\url{https://www.mathnet.ru/php/archive.phtml?jrnid=im&option_lang=rus&paperid=9597&wshow=paper}.
\textit{Formulation source.}
\end{itemize}
\end{document}
